\section{JSON-RPC}

\begin{frame}[fragile]{What Is JSON-RPC in Odoo?}
	\begin{itemize}
		\item \textbf{JSON-RPC} is Odoo's JSON-based remote call protocol.
		\item Endpoint: \texttt{/jsonrpc}
		\item Same idea as XML-RPC, but with JSON payloads.
\begin{block}{Why Use It?}
	Easier for JavaScript / modern web and mobile clients.
\end{block}
	\end{itemize}
\end{frame}

\begin{frame}[fragile]{Authenticate with JSON-RPC}
\begin{lstlisting}
import requests

url = "http://localhost:8069/jsonrpc"
payload = {
	"jsonrpc": "2.0",
	"method": "call",
	"params": {
		"service": "common",
		"method": "authenticate",
		"args": ["school_db", "admin", "admin", {}],
	},
	"id": 1,
}
uid = requests.post(url, json=payload).json()["result"]
\end{lstlisting}
\end{frame}

\begin{frame}[fragile]{Search Read via JSON-RPC}
\begin{lstlisting}
payload = {
	"jsonrpc": "2.0",
	"method": "call",
	"params": {
		"service": "object",
		"method": "execute_kw",
		"args": [
			"school_db", uid, "admin",
			"student.student", "search_read",
			[[["active", "=", True]]],
			{"fields": ["name", "email"], "limit": 10},
		],
	},
	"id": 2,
}
result = requests.post(url, json=payload).json()["result"]
\end{lstlisting}
\end{frame}

\begin{frame}[fragile]{XML-RPC vs JSON-RPC}
	\begin{itemize}
		\item Both call the same Odoo ORM layer remotely.
		\item XML-RPC: XML payloads, Python stdlib support.
		\item JSON-RPC: JSON payloads, natural for JS/web stacks.
		\item For many new integrations, JSON-RPC is more convenient.
	\end{itemize}
\end{frame}

\begin{frame}{JSON-RPC Best Practices}
	\begin{itemize}
		\item Store credentials securely; never hard-code in frontend public code.
		\item Use a dedicated integration user with least privilege.
		\item Handle \texttt{error} responses from JSON-RPC carefully.
		\item Prefer HTTPS and strong passwords/API policies.
	\end{itemize}
\end{frame}
